% TOP_PROBLEM: {"id": 3700027, "title": "Belgian Chocolate constant", "category_id": 8, "category": {"id": 8, "name": "analysis", "display_name": "Analysis", "description": "Limits, continuity, calculus, and function theory.", "slug": "analysis", "order_index": 8, "created_at": "2026-07-31T15:26:25.671Z"}, "status": "partially_solved", "problem_number": "AMR-036-0027", "set_id": 13, "statusQualification": "Uses an infimum and asks about attainment, so a limiting double-point configuration is not silently counted as admissible."}
% FIRST_POSTED: 2026-10-01
% ENTRY_KIND: statement_only
% UnsolvedMath ID 3700027; source catalogue code AMR-036-0027
% !TeX program = lualatex
% Definition and sources only; this file is not an LLM solution attempt.
\documentclass[11pt,a4paper]{article}
\usepackage[margin=25mm]{geometry}
\usepackage{fontspec}
\setmainfont{lmroman10-regular.otf}[BoldFont=lmroman10-bold.otf,
 ItalicFont=lmroman10-italic.otf,BoldItalicFont=lmroman10-bolditalic.otf]
\usepackage{amsmath,amssymb,mathtools}
\usepackage{unicode-math}
\setmathfont{latinmodern-math.otf}
\AtBeginDocument{\let\setminus\smallsetminus}
\usepackage{xurl,hyperref}
\hypersetup{colorlinks=true,urlcolor=blue}
\setcounter{secnumdepth}{0}
\setlength{\parindent}{0pt}
\setlength{\parskip}{5pt}
\setlength{\emergencystretch}{3em}
\begin{document}
\section*{Belgian Chocolate constant}\label{3700027}
{\small UnsolvedMath ID 3700027 · AMR-036-0027 · Analysis · AMR Open Problem Lists}
\subsection{Definitions and mathematical statement}
For $a\in(0,1)$, let $\mathcal B_a$ consist of the
holomorphic functions $f:\mathbb D\to\mathbb C$ satisfying
the reality condition
$f(\overline z)=\overline{f(z)}$ and the following exact
value constraints:
\[
 f^{-1}(0)=\{0\},\quad f'(0)\ne0,\qquad
 f^{-1}(1)=\{ia,-ia\},\quad
 f'(ia)f'(-ia)\ne0.
\]
All inverse images here are taken within $\mathbb D$.
Thus there is one simple zero and precisely two
distinct simple $1$-points, placed symmetrically on
the imaginary axis.

Define
\[
                 a_{\mathrm B}=
                    \inf\{a\in(0,1):\mathcal B_a\ne\varnothing\}.
\]
Determine this sharp constant, and decide whether an
admissible function realizes it. If it does, characterize
the extremal functions; if it does not, describe a
sequence of admissible parameters and functions approaching
the infimum.

This is the disk formulation of the Belgian Chocolate
problem. The symmetry of the two prescribed $1$-points
does not by itself replace the separate reality condition
on the whole function. Holomorphic functions of arbitrary
growth toward the boundary are allowed, and neither
polynomial nor rational realizations are required.
The strict inequalities $0<a<1$ keep both $1$-points
distinct from each other, from the zero, and from the
boundary. Numerical constructions yield admissible upper
bounds only when their global zero and $1$-point
counts are also established.
\subsection{Short English statement}
How close to the origin can the two symmetric imaginary one-points of a real holomorphic disk function lie?
\subsection{Sources}
\begin{itemize}
\item[S1] ULAM.AI and the UnsolvedMath Contributors, supplied problems.json, record 3700027 (AMR-036-0027), AMR Open Problem Lists. Catalogue identity and source transcription; CC BY 4.0.
\url{https://huggingface.co/datasets/ulamai/UnsolvedMath}.
\item[S2] Eremenko - My favorite unsolved problems, Goldberg relatives, Problem 3. Original source indexed by AMR-036-0027.
\url{https://www.math.purdue.edu/~eremenko/uns1.html}.
\item[S3] A. Eremenko, Goldberg's constant and its relatives, Problem 3. Author-maintained primary problem note.
\url{https://www.math.purdue.edu/~eremenko/dvi/goldconst.pdf}.
\end{itemize}
\end{document}
